% ============================================================
% Methodology
% Spine: trace scoring + language invariance for multilingual
% numerical claim verification (English, Spanish, Arabic).
% Uses plain \cite{} so it is compatible with any bibliography
% backend (natbib, biblatex, or the default BibTeX). All citation
% keys resolve against mainbib.bib.
%
% FIGURES: requires \usepackage{graphicx} (and, if including the SVGs
% directly, \usepackage{svg}). One NEW diagram is still to be drawn ---
% the language-invariance objective (Figure~\ref{fig:language-invariance}) ---
% which is currently a compiling placeholder box; all other figures are
% reused from the CLEF-2026 working note.
% ============================================================

\section{Methodology}
\label{sec:methodology}

This section explains how we investigate whether language invariance can serve as a useful training
objective for multilingual reasoning. Our starting point is the supervised trace-ranking system
developed for our CLEF~CheckThat!~2026 Task~2 submission~\cite{khawaja2026gipplab,
clef-checkthat:2026:task2}, which already includes value-aware numerical components---numeric
embeddings and plain-text numerical normalization---aimed at reducing variation caused by differences
in numeral script and format, particularly in Arabic. We use this system as the experimental backbone
and extend it in two main ways.

First, we introduce a score-level language-invariance objective that encourages the model to assign
similar scores to corresponding reasoning traces across languages. This allows us to examine whether
optimizing for invariance improves ranking performance and robustness (RQ1 and RQ3). Second, we
define an evaluation protocol that measures language invariance separately from task performance. This
makes it possible to study the relationship between cross-lingual consistency and reasoning quality
rather than assuming that more consistent predictions are necessarily more accurate (RQ2). Within this
setting we also revisit the inherited value-aware numerical components as a complementary lever for
the numerical, script-sensitive dimension of the task, now studied in combination with the invariance
objective (RQ3).

The remainder of the section presents the task and dataset, followed by the trace-scoring architecture
and its training objective. We then describe the language-invariance objective and the numerical
components, introduce the baselines, and define the evaluation protocol. Finally, we provide the
experimental configuration and the experiment matrix used for Mistral-7B, Llama-3.1-8B, and Qwen3-8B.

\subsection{Task Formulation}
\label{subsec:method-task}

We adopt the verifier-guided reranking formulation of CLEF~CheckThat!~2026
Task~2~\cite{clef-checkthat:2026:task2}, which itself follows the test-time scaling setup for numerical
claim verification~\cite{chungkham2025thinkrightmoretesttime,v2024quantemprealworldopendomainbenchmark}. Each example consists of a claim $c$, a
set of retrieved evidence snippets $E(c)$, and a set of candidate reasoning traces
$R(c)=\{1,\dots,n_c\}$ that were generated by prompting a language model at several decoding
temperatures to induce diverse reasoning paths. Each trace $i\in R(c)$ carries a trace-level verdict
$v_i$, and the claim carries a gold claim-level label $y^\star_c$. The label space is
$\mathcal{Y}=\{\textsf{True},\textsf{False},\textsf{Conflicting}\}$. A system must produce two
outputs for every claim: a ranking $r(c)$ of the candidate traces by their usefulness for reaching
the correct verdict, and a predicted claim verdict $\hat{y}_c\in\mathcal{Y}$.

The official evaluation is performed independently for each language
$\ell\in\mathcal{L}=\{\text{English},\text{Spanish},\text{Arabic}\}$. Ranking quality is measured by
Recall@5 over the traces whose verdict matches the gold claim label,
\begin{equation}
\label{eq:recall5}
\mathrm{Recall@5}(c)=
\frac{\bigl|\{\,i\in r_{1:5}(c)\;:\;v_i=y^\star_c\,\}\bigr|}
     {\bigl|\{\,i\in R(c)\;:\;v_i=y^\star_c\,\}\bigr|},
\end{equation}
where $r_{1:5}(c)$ denotes the five top-ranked traces; claim-verdict quality is measured by macro-F1
over $\mathcal{Y}$; and the per-language score is the mean of macro-F1 and Recall@5. We additionally
report MRR@5 as a secondary ranking diagnostic. A structural property of
Equation~\eqref{eq:recall5} matters throughout: for a claim with $m$ gold-matching traces, even a
perfect ranking achieves at most $\min(5,m)/m$, and a claim with $m=0$ contributes Recall@5 of zero.
On the validation data this caps the attainable Recall@5 at $0.39$, $0.32$, and $0.28$ for English,
Spanish, and Arabic respectively, which bounds the headroom of every ranking method and must be kept
in mind when reading results.

Generative ranking formulations---listwise and pairwise prompting, and DPO over complete ranking
outputs---are sensitive to candidate order and prompt format rather than content
alone~\cite{qin2024largelanguagemodelseffective,shi2025judgingjudgessystematicstudy,rafailov2024directpreferenceoptimizationlanguage},
in line with the presentation-bias literature surveyed in
section~\ref{sec:related-work}~\cite{qiao2026llmbasedlistwisererankingeffect,fu2025reliablemultilingualllmasajudge}.
Supervised trace scoring, by contrast, assigns each claim--trace pair an independent relevance score
and ranks deterministically by that score, removing the dependence on output order. We therefore take
the trace scorer as the substrate on which the language-invariance objective is defined, because it
exposes a continuous per-trace score that can be aligned across languages without having to reconcile
discrete generated rankings.

\subsection{Data}
\label{subsec:method-data}

We use the multilingual numerical claim-verification dataset released for the shared task. The
English portion is based on QuanTemp~\cite{v2024quantemprealworldopendomainbenchmark}, while the
Spanish and Arabic portions contain $2{,}808$ and $3{,}260$ claims, respectively, and come from the
previous edition of the lab together with their evidence
collections~\cite{clef-checkthat:2026:task2}. Each example contains the claim text, a collection of
evidence snippets, a list of candidate reasoning traces (\texttt{Reasoning\_traces}), the verdict
produced by each trace (\texttt{Verdict\_list}), and the gold claim label (\texttt{label}). We map
both the gold labels and trace verdicts to the common label space introduced in
Section~\ref{subsec:method-task}. A trace is then marked as relevant when its normalized verdict
matches the normalized gold label.

\paragraph{Complete and positive-only splits.}
We use two versions of each training split. The \emph{complete} version contains every claim,
including claims for which none of the candidate traces produces the correct verdict. The
\emph{positive-only} version excludes these claims. This distinction is necessary because a claim
with no gold-matching trace provides no positive example for the binary trace scorer: every
candidate trace is labelled as irrelevant. Including many such claims increases the imbalance
between positive and negative traces without helping the model learn to distinguish useful reasoning
from incorrect reasoning. We therefore apply this filtering only during training. All evaluation is
performed on the complete splits so that the results reflect the full task, including claims for
which Recall@5 is necessarily zero.

\paragraph{Aligned multilingual data for invariance.}
Training with the language-invariance objective requires parallel examples across languages. In
particular, we need the same claim and candidate reasoning traces in different languages so that the
scores assigned to corresponding traces can be compared directly. We construct these aligned data
through machine translation, creating language-specific versions of each claim, its evidence, and
its reasoning traces. During translation, we preserve the original \texttt{dataset\_index} and trace
ordering. Consequently, trace $i$ for a given claim in one language is intended to correspond to
trace $i$ for the same claim in the other languages.

\subsection{Trace-Scoring Architecture}
\label{subsec:method-scorer}

\paragraph{Backbone models.}
We implement the trace scorer using three open-weight, instruction-tuned decoder models:
Mistral-7B-Instruct-v0.3~\cite{jiang2023mistral7b}, Llama-3.1-8B-Instruct~\cite{llama31instruct},
and Qwen3-8B~\cite{qwen3technicalreport}. Each backbone is loaded as a sequence-classification
model. The pretrained transformer processes the claim--trace input, and a lightweight linear head
maps its final representation to a relevance score.

Using three model families allows us to test whether the effects of the language-invariance objective
generalize beyond a single backbone. This is particularly relevant to RQ2 and RQ3 because the models
differ in their pretraining data and tokenizers, and previous work suggests that cross-lingual
consistency is related to subword-vocabulary overlap~\cite{Qi_2023,
reddy2026effectscriptsformatsllm}.

\paragraph{Scoring head.}
Given a claim--trace input $x_i$, the model produces a scalar score $s_i$ representing the estimated
relevance of trace $i$. We consider two configurations for the classification head. The default
cross-entropy head produces two logits,
$(\ell^{\text{neg}}_i,\ell^{\text{pos}}_i)$. The corresponding probability that the trace is
relevant is
\[
p_i =
\mathrm{softmax}
\left(\ell^{\text{neg}}_i,\ell^{\text{pos}}_i\right)_{\text{pos}}.
\]
For ranking, we use the logit margin
\[
s_i=\ell^{\text{pos}}_i-\ell^{\text{neg}}_i,
\]
which produces the same ordering as the positive-class probability. The alternative binary
cross-entropy head produces a single logit $\ell_i$, which is used directly as the ranking score,
$s_i=\ell_i$, with $p_i=\sigma(\ell_i)$.

Both configurations treat trace ranking as independent binary classification: each trace is scored
without reference to the position or presentation order of the other candidates. The resulting
continuous scores can then be used both to rank traces within a claim and to align corresponding
traces across languages.

\paragraph{Parameter-efficient fine-tuning.}
We fine-tune each backbone using QLoRA~\cite{dettmers2023qloraefficientfinetuningquantized,
hu2021loralowrankadaptationlarge}. The base model is loaded in 4-bit NF4 precision with a
\texttt{bf16} computation type, and its original parameters remain frozen. Trainable low-rank
adapters are added to all linear projections using
\texttt{target\_modules=all-linear}. Unless stated otherwise, the adapters use rank $r=16$, scaling
factor $\alpha=32$, dropout $0.05$, and no adapter bias.

The classification head is kept outside the quantized model and trained in full precision using
\texttt{modules\_to\_save} to \texttt{["score"]}. This is necessary because the head is initialized for the
classification task and must learn the scale of the relevance scores. We also use PyTorch
scaled dot-product attention and gradient checkpointing to reduce the memory required by long
claim--evidence--trace inputs.

This configuration makes it possible to adapt 7--8B parameter models within a modest multi-GPU
budget. It also leaves most of the pretrained multilingual representations unchanged, allowing us
to investigate whether the shared latent representations identified in previous
work~\cite{wendler2024llamasworkenglishlatent,
schut2025multilingualllmsthinkenglish} can support language-invariant trace scoring.

\paragraph{Input construction.}
Figures~\ref{fig:supervised-trace-scorer} and~\ref{fig:supervised-trace-scorer-pipeline} summarize the complete scoring process. Each claim and candidate trace
is converted into a single input with the following structure:
\begin{quote}\small\ttfamily
Claim: \{claim\}\\
Evidence snippets:\\
{}[0] \{evidence snippet\}\\
{}[1] \{evidence snippet\}\\
\dots\\
Verdict: \{trace verdict\}\\
Justification: \{cleaned trace justification\}
\end{quote}

Before constructing the input, we clean the trace justification by removing explicit
``Label:'' and ``Justification:'' markers, repeated verdict tokens, and unnecessary whitespace. The
trace verdict remains available in its dedicated field, but removing repeated label information from
the justification prevents it from dominating the reasoning text.

Input length is controlled at both the character and token levels. We limit the lengths of the claim,
individual evidence snippets, and trace justification, as well as the number of evidence snippets
included. The resulting sequence is then truncated to the model's configured maximum length. Our
compact configuration uses two evidence snippets and a maximum sequence length of $1{,}024$ tokens,
making it suitable for rapid experimentation. The long-context configuration includes up to twenty
evidence snippets, allows up to $3{,}078$ characters per snippet, and supports sequences of up to
$15{,}000$ tokens. Traces whose cleaned justifications contain fewer than three words are excluded
because they provide too little reasoning content to be useful for training.

\begin{figure}[htbp]
\centering
  \includesvg[width=\linewidth,height=0.65\textheight,
    keepaspectratio]{trace-scorer-ranking}
\caption{Input-output flow of the supervised trace scorer. Candidate traces are scored independently, sorted by positive-class probability, and aggregated through the top-ranked trace verdicts.}
\label{fig:supervised-trace-scorer}
\end{figure}

\begin{figure}[htbp]
\centering
  \includesvg[width=\linewidth,height=0.65\textheight,
    keepaspectratio]{trace-scorer-pipeline}
\caption{Input-output flow of the supervised trace scorer. Candidate traces are scored independently, sorted by positive-class probability, and aggregated through the top-ranked trace verdicts.}
\label{fig:supervised-trace-scorer-pipeline}
\end{figure}

\subsection{Trace-Classification Objective}
\label{subsec:method-objective}

The trace scorer is trained as a binary classifier. For trace $i$ belonging to claim $c$, the target
label is defined as
\begin{equation}
\label{eq:label}
z_i=
\begin{cases}
1, & \text{if } v_i=y^\star_c,\\
0, & \text{otherwise.}
\end{cases}
\end{equation}
A trace is therefore considered relevant when its verdict agrees with the gold claim label. With the
two-logit head, we minimize the standard cross-entropy loss between
$(\ell^{\text{neg}}_i,\ell^{\text{pos}}_i)$ and $z_i$. With the single-logit head, we use binary
cross-entropy between $\sigma(\ell_i)$ and $z_i$. We refer to this classification term as the
ranking loss, $\mathcal{L}_{\text{rank}}$, because it trains the scores from which the final ranking
is constructed.

At inference time, every candidate trace is scored independently. The ranking is obtained by sorting
the traces in decreasing order of score:
\[
r(c)=\operatorname{argsort}_{i\in R(c)}(-s_i).
\]
When two traces receive the same score, their original trace indices are used to break the tie
deterministically.

The claim-level verdict is derived from the ranked traces rather than predicted by a separate
classification head. Specifically, we take a majority vote over the verdicts of the top-$k$ traces,
using $k=5$ by default:
\begin{equation}
\label{eq:verdict}
\hat{y}_c=
\operatorname{majority}
\bigl(\{\,v_i \mid i\in r_{1:k}(c)\,\}\bigr).
\end{equation}
If the vote is tied, the label that appears first in the ranking is selected. The ranking and
claim-level verdict are therefore both determined by the same learned trace scores. Improving these
scores can consequently affect both Recall@5 and macro-F1, even though the final verdict is produced
through a separate majority-vote readout.

\subsection{Language-Invariance Objective}
\label{subsec:method-invariance}

The central methodological contribution of the thesis is to add, on top of
$\mathcal{L}_{\text{rank}}$, a score-level regularizer that penalizes a model when it assigns
different relevance profiles to the \emph{same} traces across languages. This operationalizes the
measurement framing of Chapter~\ref{sec:related-work}: invariance is encouraged at the level of
continuous trace scores rather than discrete verdicts, in the spirit of cross-lingual consistency
objectives~\cite{zheng2021consistencyregularizationcrosslingualfinetuning,wang2025calmunleashingcrosslingualselfaligning,liu2026posttraininglanguagemodelscrosslingual}, but---unlike that prior
work---it is tied directly to a ranking task so that its effect on ranking quality can be measured
(RQ1). Figure~\ref{fig:language-invariance} gives an overview of the objective.

\paragraph{Aligned multilingual batching.}
Training consumes the aligned data of Section~\ref{subsec:method-data}. We first build trace-scorer
features independently per language, tagging each with its claim identity (\texttt{dataset\_index}),
trace index, and a language id. We then form \emph{aligned claim groups}: for each claim that is
present in every language, we take the set of trace indices that survive preprocessing in all
languages, and emit one group containing, for that claim, every (language, common-trace) example.
Groups are the unit of batching, so that whenever the invariance term is computed the batch is
guaranteed to contain the parallel scores it needs; a claim that has no common trace across languages
contributes no invariance signal and is dropped from the aligned set (it can still contribute to
$\mathcal{L}_{\text{rank}}$). The standard classification labels and metadata are carried through the
collator unchanged, so $\mathcal{L}_{\text{rank}}$ is identical to the monolingual case.

\begin{figure}[htbp]
  \centering
  \includesvg[width=\linewidth,height=0.65\textheight,
    keepaspectratio]{invariance2}
\caption{Language-invariant multilingual trace-scoring pipeline. Parallel English, Spanish, and Arabic claim–trace examples are aligned by claim and trace indices and processed by a shared QLoRA scorer. Training combines binary trace-ranking loss with pairwise language-invariance loss computed over normalized, aligned score vectors.}
\label{fig:language-invariance}
\end{figure}

\paragraph{Score normalization.}
Absolute score scales are not comparable across languages---a model may be systematically more or
less confident in one language even when its \emph{relative} ordering of traces is identical. Because
ranking depends only on the relative scores within a claim, we compare languages after removing
per-claim, per-language location and scale. For a claim $c$, language $\ell$, and the set of common
trace indices $T_c$, let $\mathbf{s}^{(\ell)}_c=(s^{(\ell)}_{c,i})_{i\in T_c}$ be the score vector;
we standardize it as
\begin{equation}
\label{eq:znorm}
\tilde{\mathbf{s}}^{(\ell)}_c=
\frac{\mathbf{s}^{(\ell)}_c-\operatorname{mean}\bigl(\mathbf{s}^{(\ell)}_c\bigr)}
     {\operatorname{std}\bigl(\mathbf{s}^{(\ell)}_c\bigr)+\varepsilon},
\end{equation}
with a small $\varepsilon$ for numerical stability. Standardizing makes the objective penalize
differences in the \emph{shape} of the score profile---which trace is ranked above which, and by how
much relative to the spread---rather than harmless constant offsets, aligning the regularizer with
what the ranking metric actually rewards.

\paragraph{Invariance loss.}
For each aligned claim group with at least two languages, we penalize the mean squared difference
between the standardized score vectors of every language pair, and average over claims:
\begin{equation}
\label{eq:invloss}
\mathcal{L}_{\text{inv}}=
\operatorname{mean}_{c}\;
\frac{1}{\binom{|\mathcal{L}_c|}{2}}
\sum_{\ell<\ell'}
\bigl\lVert \tilde{\mathbf{s}}^{(\ell)}_c-\tilde{\mathbf{s}}^{(\ell')}_c \bigr\rVert_2^2 \big/ |T_c|,
\end{equation}
where $\mathcal{L}_c$ is the set of languages available for claim $c$. The total training objective is
\begin{equation}
\label{eq:total}
\mathcal{L}=\mathcal{L}_{\text{rank}}+\lambda_{\text{inv}}(t)\,\mathcal{L}_{\text{inv}},
\end{equation}
where the weight $\lambda_{\text{inv}}(t)$ is warmed up linearly from $0$ to its target value over the
first fraction $\rho$ of training steps and held constant thereafter,
\begin{equation}
\label{eq:warmup}
\lambda_{\text{inv}}(t)=
\begin{cases}
\lambda^\star\cdot \dfrac{t}{\rho\,T}, & t<\rho T,\\[1.2ex]
\lambda^\star, & t\ge \rho T,
\end{cases}
\end{equation}
with $T$ the total number of optimization steps. The warm-up lets the head first learn a sensible
score scale from $\mathcal{L}_{\text{rank}}$ alone, before invariance pressure is applied; applying it
from step zero risks collapsing the scores toward a trivially consistent but uninformative solution,
the cross-lingual analogue of being ``consistently wrong''~\cite{elazar2021measuringimprovingconsistencypretrained,Qi_2023}.
Setting $\lambda^\star=0$ recovers the plain multilingual trace scorer, which gives a clean ablation:
the same data, batching, and architecture, with and without the invariance term. We log
$\mathcal{L}_{\text{rank}}$, $\mathcal{L}_{\text{inv}}$, and $\lambda_{\text{inv}}(t)$ separately
throughout training so that the two pressures can be monitored independently.

% ============================================================
% NEW DIAGRAM REQUIRED: language-invariance objective.
% This figure does not yet exist (no counterpart in the CLEF working note).
% The framed box below is a PLACEHOLDER that compiles without an image so the
% space and caption are reserved. Once the diagram is drawn, replace the
% \fbox{...} block with:
%     \includegraphics[width=0.95\linewidth]{figures/language_invariance_objective}
% Suggested content for the diagram: the same claim with its candidate traces shown
% in English / Spanish / Arabic columns -> the shared trace scorer produces a score
% vector per language -> per-claim z-normalisation -> pairwise MSE between the aligned
% score vectors (L_inv) added, with warm-up weight lambda_inv, to the ranking loss L_rank.
% ============================================================

\subsection{Value-Aware Numerical Components}
\label{subsec:method-numeric}

Numerical expressions can vary considerably across languages even when they represent the same
quantity. For example, a value may be written using Arabic-Indic or Persian digits instead of
Western Arabic digits, or it may use different conventions for decimal and thousands separators.
These surface differences can produce very different tokenizations and are a known source of
cross-lingual inconsistency~\cite{reddy2026effectscriptsformatsllm,
wallace2019nlpmodelsknownumbers,
bhattacharya2025investigatinginteractionlinguisticmathematical}. To reduce this variation, we
introduce two optional value-aware components. Both can be used with the standard trace scorer and
combined with the language-invariance objective.

\paragraph{Numeric canonicalization.}
Both components use the same numeric canonicalization procedure. Numeric spans are first identified
using a regular expression that supports Western Arabic, Arabic-Indic, and Persian digits, as well
as percentage signs, per-mille signs, and numbers containing comma or decimal separators. Each
detected expression is converted to ASCII digits, and its separators are interpreted and rewritten
using a consistent decimal format. Percentage and per-mille expressions are converted to their
corresponding decimal values. The result is then parsed as a numeric value and rendered as a
canonical decimal string.

This process maps different surface forms of the same quantity to a common representation. For
example, quantities written using different digit scripts or separator conventions receive the same
canonical value. The two numerical components described below use this representation in different
ways.

\paragraph{Value-aware numeric embeddings.}
Figure~\ref{fig:numeric-embedding} illustrates the value-aware embedding pipeline. Each detected
number is preceded by a special \texttt{<num>} token while its original surface form remains in the
input. The corresponding canonical value is passed to a small learned value encoder. This encoder
represents the canonical string character by character using a fixed vocabulary containing digits,
a sign, a decimal point, and an exponent marker. The character embeddings are processed by a
single-layer GRU and projected to the hidden size of the backbone model.

At every \texttt{<num>} position, the standard token embedding is replaced by the resulting value
vector. The transformer therefore receives an explicit representation of the number's value while
still retaining its original textual form. This reduces its dependence on the subword tokenization
of the numeral and follows previous work on value-aware numerical
representations~\cite{dutulescu2026valueawarenumericalrepresentationstransformer,
schwartz2024numerologicnumberencodingenhanced}.

To support this component, the \texttt{<num>} token is added to the tokenizer and the model's input
embedding matrix is resized. The value encoder is trained jointly with the LoRA adapters and
classification head. It is also saved with the trained adapter so that the same numerical
representation can be restored during inference.

\begin{figure}[htbp]
  \centering
  \includesvg[width=\linewidth,height=0.65\textheight,
    keepaspectratio]{numerical-embedding}
\caption{Numeric embedding pipeline. Numerical spans are canonicalized, marked with \texttt{<num>}, encoded with a GRU value encoder, and inserted as value-aware embeddings before trace scoring.}
\label{fig:numeric-embedding}
\end{figure}

\paragraph{Plain-text numeric normalization.}
The second component provides a simpler way of exposing canonical values to the model. Rather than
changing the embedding layer, it appends a short plain-text section headed
\texttt{Normalized numbers:} to the trace-scorer input. This section lists the numbers found in the
claim, evidence, and reasoning trace. When a number's original and canonical forms differ, it is
shown as a \texttt{surface=canonical} pair.

This approach requires no change to the model architecture because the normalized values are
presented as ordinary input text. It provides a lightweight alternative to the learned value
encoder and helps determine whether any improvement comes simply from making canonical values
visible to the model or from the learned embedding mechanism itself. The normalization procedure
also ignores small leading integers used as list markers, such as ``1.'' or ``2)'', so that they
are not incorrectly treated as quantities.

\subsection{Baselines}
\label{subsec:method-baselines}

We compare the trace scorers against the progression of baselines established for the shared task,
which also serve as the reference points for the invariance analysis.

\begin{itemize}
\item \textbf{Prompt baseline.} A frozen instruction-tuned backbone, used without any task-specific
fine-tuning, is our primary untrained reference. Given the claim, evidence, and candidate traces, it
is prompted to produce a trace ranking and a claim verdict directly. We run this prompted baseline in
two modes, listwise and pairwise, described next.

\item \textbf{Listwise prompting} (Figure~\ref{fig:listwise-prompted-ranking}). The model receives the claim,
evidence, and all candidate traces in one prompt and is asked to emit strict JSON containing a
permutation of trace indices and one verdict. Despite an explicit instruction not to copy the input
order, the produced ranking remains sensitive to the original trace order, exhibiting the positional
bias documented for listwise rankers~\cite{qiao2026llmbasedlistwisererankingeffect,shi2025judgingjudgessystematicstudy}.

\item \textbf{Pairwise prompting} (Figure~\ref{fig:pairwise-prompted-ranking}). Ranking is reformulated as a set of two-trace comparisons whose
accumulated wins define the ordering~\cite{qin2024largelanguagemodelseffective}. To counter A/B position bias we use a
\emph{balanced} orientation, which alternates which trace is shown first, and a \emph{bidirectional}
orientation, which evaluates each pair in both orders and counts a win only when both orders agree.
These reduce but do not remove the order effect, at a cost of $n(n-1)/2$ or $n(n-1)$ model calls per
claim.

\item \textbf{Direct Preference Optimization (DPO)} (Figures~\ref{fig:dpo-pipeline} and~\ref{fig:dpo-ranking}). An English LoRA model, initialized from a
supervised fine-tuned adapter on Mistral-7B, is trained with DPO~\cite{rafailov2024directpreferenceoptimizationlanguage} on
automatically constructed preference pairs: the chosen output ranks gold-matching traces first and
predicts the gold verdict, while rejected outputs are hard negatives (a non-matching trace placed
first, swapped top traces, or an incorrect verdict). Because DPO optimizes preferences between
complete generated rankings rather than per-trace scores, it inherits the order/format sensitivity of
generative ranking and does not directly optimize Recall@5, which is precisely the limitation the
trace scorer was designed to avoid.
\end{itemize}

% Figures reused from the CLEF-2026 working note. Export each to .pdf/.png for pdflatex,
% or load the SVGs directly with \usepackage{svg} and \includesvg{...}.
% Available in the repo under figures/: listwise_prompted_ranking, pairwise_prompted_ranking,
% dpo_input_output. The DPO pipeline figure (figures/dpo_ranking_pipeline) is the multi-stage
% translation -> preference-pair -> DPO -> cross-lingual evaluation diagram from the working note.

\begin{figure}[htbp]
  \centering
  \includesvg[width=\linewidth,height=0.65\textheight,
    keepaspectratio]{listwise-ranking}
\caption{Listwise prompted ranking setup. A frozen language model receives the claim, evidence, and all candidate traces in one prompt, then generates a complete ranking and final verdict.}
\label{fig:listwise-prompted-ranking}
\end{figure}

\begin{figure}[htbp]
  \centering
  \includesvg[width=\linewidth,height=0.65\textheight,
    keepaspectratio]{pairwise-ranking}
\caption{Pairwise prompted ranking setup. A frozen language model compares two traces at a time, and accumulated pairwise wins are used to produce the final ranking.}
\label{fig:pairwise-prompted-ranking}
\end{figure}

\begin{figure}[htbp]
  \centering
  \includesvg[width=\linewidth,height=0.65\textheight,
    keepaspectratio]{dpo-ranking}
\caption{DPO ranking setup. Chosen and rejected JSON outputs are automatically constructed from the gold label and trace verdicts, then used to train the model to prefer gold-consistent rankings.}
\label{fig:dpo-ranking}
\end{figure}

\begin{figure}[htbp]
  \centering
  \includesvg[width=\linewidth,height=0.65\textheight,
    keepaspectratio]{DPO-pipeline}
\caption{DPO ranking pipeline. Translated multilingual data is converted into ranking prompts and chosen/rejected preference pairs for DPO fine-tuning, followed by cross-lingual evaluation.}
\label{fig:dpo-pipeline}
\end{figure}

\subsection{Evaluation Protocol}
\label{subsec:method-eval}

In keeping with the foundational definition of consistency, we report task accuracy and language
invariance \emph{separately}, so that the relationship between them is an empirical finding rather
than a definitional artifact~\cite{elazar2021measuringimprovingconsistencypretrained,Qi_2023} (RQ2).

\paragraph{Task metrics.}
For each language we report Recall@5 and MRR@5 over gold-matching traces, macro-F1 of the predicted
claim verdict, and the official average of macro-F1 and Recall@5; we also report the per-language
Recall@5 upper bound so that absolute numbers can be read against attainable maxima. Multilingual
runs additionally report each metric per language and as a macro-average across languages.

\paragraph{Verdict invariance.}
Treating the language-conditioned verdicts $\{\hat{y}_\ell(c)\}_{\ell\in\mathcal{L}}$ for a claim as
the judgments of several raters, we measure their agreement with: the fraction of claims on which all
languages agree (complete agreement); Fleiss' $\kappa$ across all languages~\cite{fleiss1971measuring};
pairwise Cohen's $\kappa$~\cite{cohen1960coefficient}; and Krippendorff's $\alpha$, which tolerates
missing variants~\cite{krippendorff2004content}. These are computed over the aligned groups so that
only genuinely parallel claims are compared.

\paragraph{Ranking invariance.}
As a finer diagnostic we quantify agreement among the language-conditioned trace orderings
$\{r_\ell(c)\}$ using mean pairwise Kendall's $\tau$ and Spearman's
$\rho$~\cite{kendall1938new,spearman1904proof}, computed over the trace indices common to both
orderings, together with the top-$k$ Jaccard overlap of the ranked trace sets. The same machinery
also supports cross-\emph{variant} checks---e.g.\ original-language versus translated, or the
claim/trace-language mismatched validation sets---which isolate whether an inconsistency is driven by
the language of the claim, the language of the traces, or translation noise.

\paragraph{Training-time invariance.}
Finally, the multilingual evaluator reports $\mathcal{L}_{\text{inv}}$ itself---the mean standardized
score MSE of Equation~\eqref{eq:invloss}, overall and per language pair---on held-out data. This is
the quantity the objective directly minimizes, and tracking it alongside the behavioral
$\kappa$/$\tau$ statistics lets us check whether reducing score-level disagreement actually
translates into the verdict- and ranking-level invariance we ultimately care about, and whether it
does so without sacrificing accuracy~\cite{ifergan2024beneathsurfaceconsistencyexploring}.

\subsection{Experimental Configuration}
\label{subsec:method-config}

\paragraph{Training details.}
Unless otherwise noted, trace scorers are trained for five epochs with the AdamW optimizer, learning
rate $1\!\times\!10^{-4}$, a warm-up ratio of $0.03$ for the learning-rate schedule, per-device batch
size up to $8$ with gradient accumulation, and bf16 compute. Model selection uses early stopping with
a small patience on a validation metric---positive-class F1 for single-language runs and the
macro-F1/Recall@5 mean for multilingual runs---and the best adapter is saved at the evaluation step
where that metric peaks rather than at the last step. For the language-invariance runs we set the
invariance warm-up ratio $\rho=0.2$ and sweep the target weight $\lambda^\star$, always including
$\lambda^\star=0$ as the no-invariance control. A \emph{regularized} configuration increases adapter
dropout to $0.10$ and adds weight decay $0.01$ to reduce overfitting on the long-context inputs.
To account for run-to-run variance, every reported configuration is trained and evaluated with three
random seeds, and we report the mean across the three runs. Experiment tracking and metric logging are handled with Weights and Biases.

\paragraph{Compute.}
Experiments run on NVIDIA A100 GPUs on the institutional cluster. The main trace-scorer runs use
distributed data-parallel training across two GPUs via \texttt{torchrun}, while inference and smaller
ablations use a single GPU. All LLM-based scorers use 4-bit base-model loading with QLoRA adapters and
SDPA attention, with CUDA expandable memory segments enabled to accommodate the long
claim--evidence--trace sequences.

\paragraph{Experiment matrix.}
The three backbones play different roles in the study. \textbf{Mistral-7B-Instruct-v0.3} was carried
through the full set of ranking formulations during the shared task and is retained as a reference
point rather than extended with the invariance objective. \textbf{Llama-3.1-8B-Instruct} is the primary backbone: it has completed both
the ranking experiments (prompted baselines, DPO, and the multilingual trace scorers, including the
long-context and regularized variants) and the language-invariance experiments, and serves as the
reference against which the invariance objective is evaluated. We additionally treat plain-text
numeric normalization as a first-class setting on this backbone rather than a peripheral ablation.
\textbf{Qwen3-8B} is evaluated under the matching set of conditions---the prompted baseline, the best
ranking configuration identified for Llama (the multilingual, long-context, regularized trace
scorer), the numeric-normalization setting, and the language-invariance objective. Evaluating all
three backbones under the same conditions is what allows the thesis to separate effects of the
invariance objective from idiosyncrasies of any single model, and to study how those effects vary
across English, Spanish, and the lower-resource, different-script Arabic setting (RQ2, RQ3).